% GENERATED FILE. Edit canonical lessons, then run npm run generate:books.
% canonical-source-hash: fnv1a64:364a7cf4b1df3f75
% canonical-lessons: GU-C503-rando, GU-C503-matharvu, GU-C503-gundar, GU-C503-tirad, GU-C503-cirvu

\chapter{Smooth and Straight}
\label{ch:gu-smooth-and-straight}
\hlchaptermodality{503}

\begin{chapteropening}
\textbf{By the end of this chapter:} \emph{I can say how a carpenter smooths, glues and splits wood.}

Complete the last lesson of chapter 503: I can say how a carpenter smooths, glues and splits wood.
\end{chapteropening}

\section[\texorpdfstring{\gu{રંદો} (rando)}{રંદો (rando)}]{\gu{રંદો} (rando) — a carpenter's plane}

\label{lesson:GU-C503-rando}



\begin{quote}
Before the new one: say the Gujarati for a nail, then the Gujarati for to hammer in.
\end{quote}

\subsection*{You'll want to know: \gu{રંદો}}
\textbf{\gu{રંદો}} — \emph{rando} — \textquotedblleft{}a carpenter's plane\textquotedblright{}.

\begin{quote}
\gu{સુથાર} \gu{રંદાથી} \gu{લાકડું} \gu{સીધું} \gu{કરે} \gu{છે}.
\end{quote}

\begin{quote}
\emph{suthār randāthī lākḍũ sīdhũ kare che.}
\end{quote}

\begin{quote}
The carpenter makes the wood straight with a plane.
\end{quote}

A masculine noun.

\begin{cousinweb}
\textbf{\gu{રંદો}} (\emph{rando}) is usually traced to Persian \emph{randa}, \textquotedblleft{}a carpenter's plane\textquotedblright{}. Before \textbf{\gu{થી}} it becomes \textbf{\gu{રંદાથી}}.
\end{cousinweb}

\begin{grammarlens}[title={What you've built}]
One more word for describing what you lived and felt.
\end{grammarlens}

\subsection*{Your turn}
\begin{itemize}
\raggedright
  \item \emph{Say it:} \emph{rando}
  \item \emph{Say it:} \emph{rando}, once more
  \item \emph{Say it:} \emph{suthār randāthī lākḍũ sīdhũ kare che.} — The carpenter makes the wood straight with a plane.
  \item \emph{Recall:} say the Gujarati for a nail, then the Gujarati for to hammer in, then say \emph{rando} again
\end{itemize}

\begin{tcolorbox}[breakable,colback=teal!4,colframe=teal!35!black,title={Before you move on}]
What does \gu{રંદો} mean? (\textquotedblleft{}A carpenter's plane\textquotedblright{}.) Say it once more.
\end{tcolorbox}
\section[\texorpdfstring{\gu{મઠારવું} (maṭhārvũ)}{મઠારવું (maṭhārvũ)}]{\gu{મઠારવું} (maṭhārvũ) — to smooth (a surface)}

\label{lesson:GU-C503-matharvu}



\begin{quote}
Before the new one: say the Gujarati for to hammer in, then the Gujarati for a carpenter's plane.
\end{quote}

\subsection*{You'll want to know: \gu{મઠારવું}}
\textbf{\gu{મઠારવું}} — \emph{maṭhārvũ} — \textquotedblleft{}to smooth (a surface)\textquotedblright{}.

\begin{quote}
\gu{સુથારે} \gu{રંદાથી} \gu{બારણું} \gu{મઠાર્યું}.
\end{quote}

\begin{quote}
\emph{suthāre randāthī bārṇũ maṭhāryũ.}
\end{quote}

\begin{quote}
The carpenter smoothed the door with a plane.
\end{quote}

Past, with \gu{તેણે} (he or she …): \gu{તેણે} \gu{બારણું} \gu{મઠાર્યું}, \emph{teṇe bārṇũ maṭhāryũ}.

\begin{cousinweb}
\textbf{\gu{મઠારવું}} (\emph{maṭhārvũ}) is filed by CDIAL under Sanskrit \emph{mṛṣṭá}, \textquotedblleft{}rubbed, washed, pure\textquotedblright{}, and CDIAL gives the Gujarati verb as \textquotedblleft{}to smooth\textquotedblright{} (CDIAL 10299).
\end{cousinweb}

\begin{grammarlens}[title={What you've built}]
One more word for describing what you lived and felt.
\end{grammarlens}

\subsection*{Your turn}
\begin{itemize}
\raggedright
  \item \emph{Say it:} \emph{maṭhārvũ}
  \item \emph{Say it:} \emph{maṭhārvũ}, once more
  \item \emph{Say it:} \emph{suthāre randāthī bārṇũ maṭhāryũ.} — The carpenter smoothed the door with a plane.
  \item \emph{Recall:} say the Gujarati for to hammer in, then the Gujarati for a carpenter's plane, then say \emph{maṭhārvũ} again
\end{itemize}

\begin{tcolorbox}[breakable,colback=teal!4,colframe=teal!35!black,title={Before you move on}]
What does \gu{મઠારવું} mean? (\textquotedblleft{}To smooth (a surface)\textquotedblright{}.) Say it once more.
\end{tcolorbox}
\section[\texorpdfstring{\gu{ગુંદર} (gundar)}{ગુંદર (gundar)}]{\gu{ગુંદર} (gundar) — glue}

\label{lesson:GU-C503-gundar}



\begin{quote}
Before the new one: say the Gujarati for a carpenter's plane, then the Gujarati for to smooth.
\end{quote}

\subsection*{You'll want to know: \gu{ગુંદર}}
\textbf{\gu{ગુંદર}} — \emph{gundar} — \textquotedblleft{}glue\textquotedblright{}.

\begin{quote}
\gu{આ} \gu{ગુંદર} \gu{લાકડા} \gu{પર} \gu{સારું} \gu{ચોંટે} \gu{છે}.
\end{quote}

\begin{quote}
\emph{ā gundar lākḍā par sārũ cõṭe che.}
\end{quote}

\begin{quote}
This glue sticks well to wood.
\end{quote}

A masculine noun.

\begin{cousinweb}
\textbf{\gu{ગુંદર}} (\emph{gundar}) is filed by CDIAL under Sanskrit \emph{gundra}, the name of a reed, and CDIAL lists the Gujarati word with the sense \textquotedblleft{}gum\textquotedblright{} (CDIAL 4199); how the reed's name came to mean gum is not explained there.
\end{cousinweb}

\begin{grammarlens}[title={What you've built}]
One more word for describing what you lived and felt.
\end{grammarlens}

\subsection*{Your turn}
\begin{itemize}
\raggedright
  \item \emph{Say it:} \emph{gundar}
  \item \emph{Say it:} \emph{gundar}, once more
  \item \emph{Say it:} \emph{ā gundar lākḍā par sārũ cõṭe che.} — This glue sticks well to wood.
  \item \emph{Recall:} say the Gujarati for a carpenter's plane, then the Gujarati for to smooth, then say \emph{gundar} again
\end{itemize}

\begin{tcolorbox}[breakable,colback=teal!4,colframe=teal!35!black,title={Before you move on}]
What does \gu{ગુંદર} mean? (\textquotedblleft{}Glue\textquotedblright{}.) Say it once more.
\end{tcolorbox}
\section[\texorpdfstring{\gu{તિરાડ} (tirāḍ)}{તિરાડ (tirāḍ)}]{\gu{તિરાડ} (tirāḍ) — a crack}

\label{lesson:GU-C503-tirad}



\begin{quote}
Before the new one: say the Gujarati for to smooth, then the Gujarati for glue.
\end{quote}

\subsection*{You'll want to know: \gu{તિરાડ}}
\textbf{\gu{તિરાડ}} — \emph{tirāḍ} — \textquotedblleft{}a crack\textquotedblright{}.

\begin{quote}
\gu{જૂના} \gu{બારણામાં} \gu{એક} \gu{મોટી} \gu{તિરાડ} \gu{છે}.
\end{quote}

\begin{quote}
\emph{jūnā bārṇāmā̃ ek moṭī tirāḍ che.}
\end{quote}

\begin{quote}
There is a big crack in the old door.
\end{quote}

A feminine noun.

\begin{cousinweb}
\textbf{\gu{તિરાડ}} (\emph{tirāḍ}) is perhaps related to \emph{taraṛ}, \textquotedblleft{}a crack\textquotedblright{}, which CDIAL files under \emph{traṭ} (CDIAL 5988), but CDIAL does not list this form, so its history is not settled.
\end{cousinweb}

\begin{grammarlens}[title={What you've built}]
One more word for describing what you lived and felt.
\end{grammarlens}

\subsection*{Your turn}
\begin{itemize}
\raggedright
  \item \emph{Say it:} \emph{tirāḍ}
  \item \emph{Say it:} \emph{tirāḍ}, once more
  \item \emph{Say it:} \emph{jūnā bārṇāmā̃ ek moṭī tirāḍ che.} — There is a big crack in the old door.
  \item \emph{Recall:} say the Gujarati for to smooth, then the Gujarati for glue, then say \emph{tirāḍ} again
\end{itemize}

\begin{tcolorbox}[breakable,colback=teal!4,colframe=teal!35!black,title={Before you move on}]
What does \gu{તિરાડ} mean? (\textquotedblleft{}A crack\textquotedblright{}.) Say it once more.
\end{tcolorbox}
\section[\texorpdfstring{\gu{ચીરવું} (cīrvũ)}{ચીરવું (cīrvũ)}]{\gu{ચીરવું} (cīrvũ) — to split (wood)}

\label{lesson:GU-C503-cirvu}



\begin{quote}
Before the new one: say the Gujarati for glue, then the Gujarati for a crack.
\end{quote}

\subsection*{You'll want to know: \gu{ચીરવું}}
\textbf{\gu{ચીરવું}} — \emph{cīrvũ} — \textquotedblleft{}to split (wood)\textquotedblright{}.

\begin{quote}
\gu{અરુણે} \gu{મોટું} \gu{લાકડું} \gu{બે} \gu{ભાગમાં} \gu{ચીર્યું}.
\end{quote}

\begin{quote}
\emph{Aruṇe moṭũ lākḍũ be bhāgmā̃ cīryũ.}
\end{quote}

\begin{quote}
Arun split the big piece of wood in two.
\end{quote}

Past, with \gu{તેણે} (he or she …): \gu{તેણે} \gu{લાકડું} \gu{ચીર્યું}, \emph{teṇe lākḍũ cīryũ}.

\begin{cousinweb}
\textbf{\gu{ચીરવું}} (\emph{cīrvũ}) comes from \emph{*cīrayati}, \textquotedblleft{}splits\textquotedblright{}, a form CDIAL reconstructs, and CDIAL lists the Gujarati verb, \textquotedblleft{}to slit, crack, cut\textquotedblright{} (CDIAL 4844).
\end{cousinweb}

\begin{grammarlens}[title={What you've built}]
One more word for describing what you lived and felt.
\end{grammarlens}

\subsection*{Your turn}
\begin{itemize}
\raggedright
  \item \emph{Say it:} \emph{cīrvũ}
  \item \emph{Say it:} \emph{cīrvũ}, once more
  \item \emph{Say it:} \emph{Aruṇe moṭũ lākḍũ be bhāgmā̃ cīryũ.} — Arun split the big piece of wood in two.
  \item \emph{Recall:} say the Gujarati for glue, then the Gujarati for a crack, then say \emph{cīrvũ} again
\end{itemize}

\begin{tcolorbox}[breakable,colback=teal!4,colframe=teal!35!black,title={Before you move on}]
What does \gu{ચીરવું} mean? (\textquotedblleft{}To split (wood)\textquotedblright{}.) Say it once more.
\end{tcolorbox}
